\documentclass[10pt,a4paper]{amsart}
\usepackage[margin=19mm]{geometry}
\usepackage{amsmath,amssymb,mathtools}
\usepackage[scr=rsfs,cal=euler]{mathalfa}
\usepackage{xfrac}
\usepackage[unicode,hidelinks,bookmarks=false]{hyperref}
\newcommand{\ie}{i.e.\ }
\newcommand{\cf}{cf.\ }
\newcommand{\resp}{respectively\ }
\newcommand{\Subsection}{\subsection}
\providecommand{\nobreakdash}{\nobreak\hbox{-}}
\setlength{\emergencystretch}{2em}
\multlinegap0pt
\newcommand{\C}{\mathbb{C}}
\usepackage{amssymb}
\usepackage[shortlabels]{enumitem}
\newtheorem{lem}{Lemma}
\newcommand{\D}{\mathbb{D}}
\newcommand{\h}{{h}}
\newcommand{\kk}{{k}}
\newcommand{\my}{\mathfrak m}
\title{Zariski equisingularity of surface singularities in $\C^3$ by a local invariant}
\author{Adam Parusi\'nski, Lauren\c{t}iu P\u{a}unescu}
\date{Source: arXiv:2602.15192v1 (CC BY 4.0)}
\begin{document}
\maketitle

\noindent\emph{Source and licence.} Zariski equisingularity of surface singularities in $\C^3$ by a local invariant. Version arXiv:2602.15192v1 (CC BY 4.0), distributed by arXiv under the Creative Commons Attribution 4.0 International licence, http://creativecommons.org/licenses/by/4.0/, as recorded in the arXiv licence metadata for this exact version and shown on the abstract page https://arxiv.org/abs/2602.15192v1. Full licence notice: \texttt{licenses/arxiv-2602.15192-CC-BY-4.0.txt}.\par\smallskip

\noindent\textit{Editorial scope.} Counted unit: the article's lem lem:rational (source0.tex lines 1314--1324). The article's complete original proof of it is delivered with it (source0.tex lines 1326--1381, one \texttt{proof} environment of 56 lines). No mathematics is written, repaired or re-derived: the statement and the proof are the article's own lines, in the article's own notation, and no retained line is substituted or re-flowed.  No further source-local context is needed: every cross-reference of every retained line resolves inside the two delivered spans.  Nothing was omitted from any retained span, so the package carries no omission record.  No retained line contains a citation, so the wrapper carries no bibliography and no bibliography entry.  No retained line contains a figure environment, an image-inclusion command or drawing code, so no image is delivered and the package declares no asset dependency.  Editorial formatting only, in addition: an AMS \texttt{amsart} preamble that restates the article's own declarations and macros used by the retained lines, namely the environments \texttt{lem} on separate counters, together with the standard packages those lines need.  The retained environments print in this document as Lemma 1 in order of appearance, whereas the article prints the same items with the numbering of its own full document.  The complete native source of the article as distributed in the arXiv e-print for this exact version is bundled unchanged as \texttt{sources/194781/source0.tex}.
\begin{lem}\label{lem:rational}
Suppose that $f\in \kk(\tau) [[x]]$ is uniformly rational in $\tau$ and is regular in $x_{r+1}$ (i.e. $f(0,x_{r+1})\not\equiv 0$).  
Then, the coefficients $a_i$  
of the Weierstrass associated polynomial of $f$ over $\kk(\tau)$
\begin{align*}
f(x,\tau) = \Big (x_{r+1}^d + \sum_{i=1}^d a_i(x', \tau) x_{r+1}^{d-i} \Big ) \ unit(x,\tau), 
\end{align*}
$x'= (x_1, \ldots , x_{r})$, are uniformly rational in 
$\tau$ (maybe with a different  denominator $c$). In particular, the discriminant 
$\D_f(x',\tau)$  of   $f$ also is  uniformly rational in $\tau$.
\end{lem}

\begin{proof}
Let us consider first  $f = \sum c_n(y) z^n\in \kk[\tau] [[y,z]]$, where $y=(y_1, \ldots , y_r)$ and $z$ is a single variable, and suppose, moreover, that there is $d\ge 1$  such that $c_i(0)=0$ for $i<d $ and $c_d(0)=1 $.
We denote by $\my_y$ the ideal of $\kk (\tau) [[y,z]]$ generated by $y_1, \ldots, y_r$ and by $\my_z$ the one generated by 
$z$.  Thus $\my := \my_y + \my_z$ is the maximal ideal of $\kk (\tau) [[y,z]]$.  
 
For a series $f$ we define $\h (f):= \sum_{n\geq 0} c_{n+d} z^n$.  
Then  $\h (f)\in 1+ \my \cap \kk [\tau ] [[y,z]]$ and therefore $\h (f)^{-1}\in 1+ \my \cap \kk [\tau ] [[y,z]]$.   Write 
$$
f= z^d \h (f) + R,  \quad \text {with} \quad R\in \my_y \cap \kk[\tau] [[y]][z]_{d-1},  
$$
that is $R$ is a polynomial in $z$ of degree at most $d-1$.  
  Define 
$f_1:= (\h (f))^{-1} f = z^d +   (\h (f))^{-1} R  $. Note that 
$(\h (f))^{-1} R \in  \my_y \cap \kk [\tau ] [[y,z]]$  and 
$$
f_1= z^d \h (f_1) + R_1,  \quad R_1\in  \my_y \cap \kk[\tau] [[y]][z]_{d-1}, \quad \h (f_1)
\in 1+ \my_y \cap \kk [\tau ] [[y,z]] .$$
The next step is similar. Define 
$f_2:= (\h (f_1))^{-1} f_1 = z^d +  (\h (f_1))^{-1} R_1 $.  Then 
$$
f_2= z^d \h (f_2) + R_1 + R_2 ,  \quad R_2\in \my_y^2 \cap \kk[\tau] [[y]][z]_{d-1}, \quad \h (f_2) \in 1+ 
\my_y^2 \cap \kk [\tau ] [[y,z]] .
$$
Recursively we set $f_{k}:= (\h (f_{k-1}))^{-1} f_{k-1} = z^d + (\h (f_{k-1}) )^{-1} (\sum_{i=1}^{k-1} R_i) ,$ 
$$
f_{k} = z^d \h (f_{k}) + \sum_{i=1}^{k} R_i, \quad 
R_{k}\in  \my_y^{k} \cap \kk[\tau] [[y]][z]_{d-1}, \quad \h (f_{k})\in 1+ 
\my_y^{k} \cap \kk [\tau ] [[y,z]] .
$$
  With $k\to \infty$ we get 
\begin{align*}
f = (\prod_{i=0}^{\infty} \h (f_i)) (z^d  + \sum_{i=1}^{\infty} R_i)  \quad \text {with} \quad \sum_{i=1}^{\infty} R_i \in 
\kk[\tau] [[y]][z]_{d-1}
\end{align*}
This shows the lemma under the assumptions $f \in \kk[\tau] [[y,z]]$ and $c_d(0) =1$.  


In the general case  $f \in \kk[\tau] [[\frac y c, \frac z c]]$, $c_i(0)=0$ for $i<d $ and $c_d(0)= \frac b {c^ d} \ne 0 $.  
Set $\tilde y = \frac y {cb^{d+1}}$, $\tilde z = \frac z {cb}$.  Then $\tilde f \in \kk[\tau] [[\tilde y, \tilde z]]$, 
given by $\tilde f (\tilde y, \tilde z) := b ^ {-d-1} f(cb^{d+1}\tilde y, cb \tilde z)$ satisfies the assumptions of the case 
considered above and therefore 
\begin{align*}
\tilde f (\tilde y, \tilde z)   = \tilde u (\tilde y, \tilde z) (\tilde z^d  + \sum_{i=1}^{d} \tilde a_i (\tilde y) \tilde z^ {d-i})    \quad \text {with} \quad \tilde a_i \in 
\kk[\tau] [[\tilde y]]. 
\end{align*} 
This implies that 
\begin{align*}
 f (y, z)   & = b^{d+1}\tilde u \big (\frac y {cb^{d+1}}, \frac z {cb} \big ) \Big ( \big(\frac z {cb}\big )^d  + \sum_{i=1}^{d} \tilde a_i \big (\frac y {cb^{d+1}}\big ) \big (\frac z {cb}\big )^ {d-i} \Big)\\
 & = u(y,z) \Big (z^d  + \sum_{i=1}^{d} (cb)^i \tilde a_i \big (\frac y {cb^{d+1}}\big ) z ^ {d-i} \Big ), 
\end{align*} 
where $u(z,y) = \frac b {c^d}  \tilde u (\frac y {cb^{d+1}}, \frac z {cb}) $.  This completes the proof of lemma.    
%  Write 
% $$
% f= z^d \h (f) + R,  \quad \text {with} \quad R\in \my_y \cap \kk[\tau] [[y]][z]_{d-1}.   
% $$ 
\end{proof}

\end{document}
